% ======================================================================
% SEÇÃO 5 — DISCUSSÃO
% ======================================================================
\section{Discuss\~ao}
\label{sec:disc}

Os resultados indicam que a implementa\c{c}\~ao atende aos objetivos funcionais e
de seguran\c{c}a definidos, com 67 testes aprovados distribu\'idos pelas tr\^es
camadas do sistema. Tr\^es observa\c{c}\~oes merecem destaque. Primeiro, os testes
de FSM, PID e rampa validam matematicamente a estrat\'egia de controle mista,
incluindo o tratamento da satura\c{c}\~ao e a sele\c{c}\~ao da regi\~ao de
opera\c{c}\~ao --- condi\c{c}\~ao necess\'aria, ainda que n\~ao suficiente, para a
qualidade da rampa em bancada. Segundo, a valida\c{c}\~ao dos watchdogs em ambas
as camadas demonstra que o requisito de parada segura autom\'atica \'e atendido no
n\'ivel de software, independentemente do hardware. Terceiro, os testes de ciclo
completo com simulador confirmam a coer\^encia entre a FSM, o protocolo serial e a
telemetria consumida pela IHM, reduzindo o risco de falhas de integra\c{c}\~ao
quando o hardware real for conectado.

Dos crit\'erios de aceita\c{c}\~ao consolidados na Tabela~\ref{tab:criterios},
tr\^es foram verific\'aveis no ambiente simulado (comunica\c{c}\~ao a 4~Hz,
persist\^encia e watchdog) e dois dependem de medi\c{c}\~ao em bancada
(linearidade da rampa e estabilidade do Forno 2). Esses \'ultimos s\~ao
crit\'erios de \emph{desempenho t\'ermico} do processo f\'isico e n\~ao podem ser
atestados por testes de software: exigem a calibra\c{c}\~ao da curva de
aquecimento do Tubo U (Taxa $\times$ PWM) e a verifica\c{c}\~ao experimental das
curvas de temperatura, etapas de continuidade deste trabalho.

A principal limita\c{c}\~ao \'e, portanto, a aus\^encia de medi\c{c}\~oes em
bancada com o processo f\'isico. Al\'em disso, a raz\~ao de taxas empregada na
regi\~ao de malha aberta depende da calibra\c{c}\~ao da taxa de aquecimento do
sistema; na aus\^encia desse dado, a implementa\c{c}\~ao adota uma aproxima\c{c}\~ao
que precisa ser confirmada. Registra-se, ainda, a indefini\c{c}\~ao quanto ao
circuito integrado exato de convers\~ao do termopar (MAX6675 ou MAX31855), que
n\~ao altera a l\'ogica de controle, mas deve ser confirmada na documenta\c{c}\~ao
f\'isica do equipamento.

No plano t\'ecnico, os resultados demonstram a viabilidade da arquitetura
proposta para substituir o sistema legado, com vantagens de custo, flexibilidade
e seguran\c{c}a. No plano operacional, a persist\^encia dos par\^ametros e o
registro de tend\^encias criam condi\c{c}\~oes para a repetibilidade e a auditoria
do m\'etodo, requisitos valorizados em ambiente anal\'itico. A etapa de
valida\c{c}\~ao em bancada \'e o passo natural para confirmar os crit\'erios de
desempenho t\'ermico e calibrar a rampa do Tubo U.
